

\begin{figure}[t]
\begin{pcvstack}[boxed,center,space=1em]
    \centering
    \procedure{Algorithm $\genfork^\cdv(X)$  \fbox{$\modgenfork^\cdv(X)$}}{
    \text{Sample coins $\rho$ for $\cdv$ at random.} \\
    h_1, \ldots, h_\forkhashes  \rgets H \\
    (j, \auxstring) \rgets \cdv\bigl(X, (h_1, \ldots, h_\forkhashes); \rho\bigr) \\
    \pcreturn \bot\ \pcif j = \bot \\
    h_j', \ldots, h_\forkhashes' \rgets H \\
    (j', \auxstring') \rgets \cdv\bigl(X,  (h_1, \ldots, h_{j-1}, h_j', \ldots, h_\forkhashes'); \rho\bigr) \\
    \pcreturn \bot\ \pcif j' = \bot \\
    \pcreturn \bot\ \pcif j \neq j'  \\
    \fbox{\parbox{25em}{%
    $\text{\algcom Fail if any outputs sampled from $H$ collide} $ \\
    $\pcif (h_1, \ldots, h_\forkhashes , h_j', \ldots, h_\forkhashes')$ contains a repeated element  \\
    $\pcind \pcreturn \bot$
    }} \\
    \pcreturn (h_j, h_j', \auxstring, \auxstring')
  }
\end{pcvstack}
    \caption{
    The general forking algorithm $\genfork^\cdv(X)$ and the modified general forking algorithm $\modgenfork^\cdv(X)$.
    }
    \label{fg:gen-fork-mod}
\end{figure}




\section{General Forking Lemma}
\label{mod-fork}



We review the general forking lemma by Bellare and Neven~\cite{BellareN06},
which itself is a formalization of the forking lemma  introduced by Pointcheval and Stern~\cite{PointchevalS96}.
The forking lemma is a useful tool to proving the security of signature schemes in the random oracle model (ROM).
We use the forking lemma to prove the security of FROST in the ROM in Chapter~\ref{chp:frost},
and Storm in Chapter~\ref{chp:storm}.

We show the general forking algorithm in Figure~\ref{fg:gen-fork-mod},
as well as a slight modification,
which we describe in Section~\ref{sec:modified-forking}.

The algorithms are defined with respect to an algorithm $\cdv$ and instance $X$.
The difference between the general forking algorithm and modified variant is shown in a box, for emphasis.
In summary,
the modified general forking algorithm aborts if \emph{any} of the $(h_1, \ldots, h_\forkhashes, h_j', \ldots, h_\forkhashes')$ collide.
The modified variant allows us to account for hash collisions in a generalized manner,
as the forking algorithm aborts on any hash collision.




\begin{lemma}[General Forking Lemma~\cite{BellareN06}] \label{lem:forking-lemma}
Let $H$ be a finite set and $\forkhashes \geq 1$ be an integer.
Let $\IG$ be a randomized instance generator and let $X \rgets \IG$ be an instance.
Let $\cdv$ be a randomized algorithm that takes as input $X$,
quantities $h_1,\ldots,h_\forkhashes \in H$,
and a random tape $\rho$.
Let $\accept(\cdv)$ be the probability that $\cdv$ outputs an accepting answer, namely
\[
%\accept(\cdv) \deq \Pr_{X \rgets \IG, \{h_1, \ldots, h_q \} \rgets H} \bigl[ (j, \auxstring) \rgets \cdv(X, \{h_1, \ldots, h_q \}; \rho)\ : j \neq \bot \bigr].
\accept(\cdv) \deq
        \Pr\left[ j \neq \bot \ :\  \
               \begin{array}{c}
                               X \rgets \IG, \ \
                               h_1, \ldots, h_\forkhashes  \rgets H  \\
                               (j, \auxstring) \rgets \cdv\bigl(X, (h_1, \ldots, h_\forkhashes); \rho\bigr)
                \end{array}
            \right].
\]
Let $\genfork^\cdv(X)$ be the general forking algorithm defined in Figure~\ref{fg:gen-fork-mod} and let
\[
\accept(\genfork^\cdv) \deq
        \Pr\left[ \alpha \neq \bot \ :\  \   X \rgets \IG,\ \ \alpha \rgets \genfork^\cdv(X)  \right].
\]
Then,  $\accept(\genfork^\cdv)$ is bounded by
\[
    \accept(\genfork^\cdv)  \geq \accept(\cdv) \cdot \left(\frac{\accept(\cdv) }{\forkhashes} -\frac{1}{\abs{H}} \right).
\]
\end{lemma}


\subsubsection{A modified forking lemma.}
\label{sec:modified-forking}
We will employ a slight modification of the forking experiment,
and give a corollary on how it impacts the accepting probability of its output.

The modification to the forking experiment is natural;
intuitively,
we add an additional abort condition if any of the values $h_1, \ldots, h_\forkhashes, h_j', \ldots, h_\forkhashes'$ collide.
Because there are at most $2\forkhashes$ values, and they are all sampled uniformly at random from the set $H$,
the probability that any of them collide is at most $(2\forkhashes)^2/\abs{H}$.
By considering this case here,
we can be sure that such collisions are considered in our proof of security for \glitter.

\begin{corollary} \label{cor:modified-fork}
    Let $\modgenfork^\cdv$ be the forking experiment in Figure~\ref{fg:gen-fork-mod}.
    Then using the notation of Lemma~\ref{lem:forking-lemma} we have
    \begin{align} \label{eq:bound-mod-fork}
    \begin{split}
        \accept(\modgenfork^\cdv) & \geq \accept(\genfork^\cdv) - \Pr[\badhashevent]   \\
        & \geq \accept(\genfork^\cdv) - (2\forkhashes)^2/\lvert H \rvert
    \end{split}
    \end{align}
    Where $\badhashevent$ denotes the event that $\modgenfork^\cdv$ returns $\bot$ due to the boxed lines in Figure~\ref{fg:gen-fork-mod}.
\end{corollary}

\paragraph{Proof Intuition for Corollary~\ref{cor:modified-fork}.}
The intuition for Corollary~\ref{cor:modified-fork} is straightforward,
as it simply counts the number of times two values sampled uniformly at random from a $H$ (with replacement) might collide.
As such,
it counts the number of times the same value might be sampled twice,
when $\forkhashes $ total sampling events are performed.


